\section{Completeness and compactness}
\label{sec:prop-metatheory-03}

\begin{theorem}[Strong completeness of classical propositional natural deduction]
\label{thm:prop-completeness}
For every \(\Gamma\subseteq\Fm_{\CPL}\) and every
\(\varphi\in\Fm_{\CPL}\),
\[
  \Gamma\entails\varphi
  \quad\Longrightarrow\quad
  \Gamma\derives_{\ND_{\CPL}}\varphi.
\]
\end{theorem}
\begin{proof}
We prove the contrapositive.  Suppose
\(\Gamma\notderives_{\ND_{\CPL}}\varphi\).  Then
\(\Gamma\cup\set{\neg\varphi}\) is consistent.  For suppose a derivation
\(D\) of \(\bottom\) has open assumptions in
\(\Gamma\cup\set{\neg\varphi}\).  If no open occurrence of
\(\neg\varphi\) appeared in \(D\), then \(D\) would witness
\(\Gamma\derives_{\ND_{\CPL}}\bottom\), and \(\bottom E\) would give
\(\Gamma\derives_{\ND_{\CPL}}\varphi\), contrary to the supposition.  So
at least one open occurrence remains; mark all open occurrences of
\(\neg\varphi\) in \(D\) with a common fresh label \(i\), by the merging
relabeling of the derivation definition, and apply \(\mathsf{RAA}\) with
label \(i\) at the root, discharging every such occurrence, to obtain
\(\Gamma\derives_{\ND_{\CPL}}\varphi\)---again contrary to the
supposition.  Hence \(\Gamma\cup\set{\neg\varphi}\) is consistent.  By
\cref{lem:prop-lindenbaum}, extend this set to a maximal
consistent set \(\Delta\).  The truth lemma gives
\(v_\Delta\entails\Gamma\) and
\(v_\Delta\entails\neg\varphi\).  Hence
\(v_\Delta\notentails\varphi\), so
\(\Gamma\notentails\varphi\).
\end{proof}

Combining \cref{thm:prop-nd-soundness,thm:prop-completeness} yields
\[
  \Gamma\derives_{\ND_{\CPL}}\varphi
  \quad\text{iff}\quad
  \Gamma\entails\varphi.
\]

\begin{theorem}[Propositional compactness]
\label{thm:prop-compactness}
For every \(\Gamma\subseteq\Fm_{\CPL}\) and every
\(\varphi\in\Fm_{\CPL}\):
\begin{enumerate}
  \item if \(\Gamma\entails\varphi\), then some finite
        \(\Gamma_0\subseteq\Gamma\) satisfies
        \(\Gamma_0\entails\varphi\); and
  \item \(\Gamma\) is satisfiable iff every finite subset of \(\Gamma\) is
        satisfiable.
\end{enumerate}
\end{theorem}
\begin{proof}
If \(\Gamma\entails\varphi\), completeness gives a finite derivation from
\(\Gamma\).  Only finitely many members of \(\Gamma\) occur as open
assumptions; call their set \(\Gamma_0\).  Soundness gives
\(\Gamma_0\entails\varphi\), proving 1.

The forward implication in 2 is immediate.  For the converse, if \(\Gamma\)
were not satisfiable, then \(\Gamma\entails\bottom\).  Clause 1 would give a
finite \(\Gamma_0\subseteq\Gamma\) with
\(\Gamma_0\entails\bottom\), so \(\Gamma_0\) would not be satisfiable.
\end{proof}

\begin{theorem}[Completeness of the sequent calculus]
\label{thm:prop-lk-completeness}
For all finite multisets \(\Gamma,\Delta\) of propositional formulas, the
sequent \(\Gamma\Rightarrow\Delta\) is derivable in \(\LK_{\CPL}\) iff
every valuation making every formula of \(\operatorname{supp}(\Gamma)\)
true makes some formula of \(\operatorname{supp}(\Delta)\) true.
\textit{Status: proof sketch.}
\end{theorem}
\begin{proof}[Proof sketch]
The exact scope is the displayed cut-admitting calculus with finite multiset
contexts.  Left to right is \cref{thm:prop-lk-soundness}, whose semantic
condition on a sequent is exactly the one displayed here.

For right to left, first record the translation subroutine.  Induction on
a \(\ND_{\CPL}\)-derivation translates every introduction rule to the
corresponding right rule, every elimination rule to the corresponding
left rule followed when necessary by cut, and assumption discharge to a
right rule.  Weakening and contraction align the finite contexts.  The
induction establishes that if
\(\Gamma_0\derives_{\ND_{\CPL}}\varphi\) for a finite set \(\Gamma_0\),
then the sequent \(\Gamma_0\Rightarrow\varphi\) is derivable in the
displayed calculus.  The rule-by-rule translation trees, with the exact
count of contractions and weakenings, are
\cref{def:p02-w02b-p02-009-nd-to-lk}, proved correct in
\cref{thm:p02-w02b-p02-009-nd-to-lk}; the translation uses Cut and does not
establish cut elimination, which is
\cref{thm:p02-w02b-p02-009-cut-admissibility}.

Now suppose every valuation making every formula of
\(\operatorname{supp}(\Gamma)\) true makes some formula of
\(\operatorname{supp}(\Delta)\) true, and list
\(\operatorname{supp}(\Delta)=\set{\delta_1,\ldots,\delta_m}\) with
\(m\geq0\).
\begin{enumerate}
  \item Let \(S\) be the set of members of
        \(\operatorname{supp}(\Gamma)\) together with
        \(\neg\delta_1,\ldots,\neg\delta_m\).  A valuation satisfying
        all of \(S\) would make every member of
        \(\operatorname{supp}(\Gamma)\) true, hence some
        \(\delta_j\) true together with \(\neg\delta_j\).  So no
        valuation satisfies \(S\), and \(S\entails\bottom\).
  \item \Cref{thm:prop-completeness} gives
        \(S\derives_{\ND_{\CPL}}\bottom\).  The open assumptions of
        that derivation form a finite subset of \(S\), say
        \(\gamma_1,\ldots,\gamma_k\in\operatorname{supp}(\Gamma)\) and
        \(\neg\delta_{j_1},\ldots,\neg\delta_{j_r}\).  The translation
        subroutine turns it into an \(\LK_{\CPL}\)-derivation of
        \(\gamma_1,\ldots,\gamma_k,
        \neg\delta_{j_1},\ldots,\neg\delta_{j_r}\Rightarrow\bottom\).
  \item Cut with the initial \(\bottom\)-sequent
        \(\bottom\Rightarrow\), the instance of
        \(\Gamma,\bottom\Rightarrow\Delta\) with both contexts empty,
        empties the succedent.
  \item For each \(j\), the sequent
        \(\Rightarrow\delta_j,\neg\delta_j\) is derivable: start from
        the identity sequent \(\delta_j\Rightarrow\delta_j\) and apply
        \(\neg R\).  Cutting it against the current sequent, with cut
        formula \(\neg\delta_j\),
        \[
         \Rule[Cut]{\Rightarrow\delta_j,\neg\delta_j\qquad
          \Gamma',\neg\delta_j\Rightarrow\Pi}
          {\Gamma'\Rightarrow\delta_j,\Pi},
        \]
        moves \(\delta_j\) into the succedent; after \(r\) such cuts
        the derived sequent is
        \(\gamma_1,\ldots,\gamma_k\Rightarrow
        \delta_{j_1},\ldots,\delta_{j_r}\).
  \item Weakening (\(W_L\), \(W_R\)) adds the unused members of
        \(\Gamma\) and \(\Delta\), and contraction (\(C_L\),
        \(C_R\)) merges duplicates, giving exactly the multiset
        sequent \(\Gamma\Rightarrow\Delta\).
\end{enumerate}
When \(m=0\), the hypothesis asserts that no valuation makes every
member of \(\operatorname{supp}(\Gamma)\) true, step 4 performs no cuts,
and the same route derives \(\Gamma\Rightarrow\Delta\) with empty
succedent; an empty \(\Gamma\) is likewise covered.  The route uses cut,
which is primitive in the displayed calculus, and makes no cut-free
claim; no choice principle enters, because the propositional Lindenbaum
lemma runs on a countable enumeration.  The status is a sketch because
the individual translation trees of the subroutine remain omitted.
\end{proof}

